%%%%%%%%%%%%%%%%%%%%%%%%%%%%%%%%%%%%%%%%%
% Medium Length Professional CV
% LaTeX Template
% Version 2.0 (8/5/13)
%
% This template has been downloaded from:
% http://www.LaTeXTemplates.com
%
% Original author:
% Trey Hunner (http://www.treyhunner.com/)
%
% Important note:
% This template requires the resume.cls file to be in the same directory as the
% .tex file. The resume.cls file provides the resume style used for structuring the
% document.
%
%%%%%%%%%%%%%%%%%%%%%%%%%%%%%%%%%%%%%%%%%

%----------------------------------------------------------------------------------------
%	PACKAGES AND OTHER DOCUMENT CONFIGURATIONS
%----------------------------------------------------------------------------------------

\documentclass{resume} % Use the custom resume.cls style

\usepackage[left=0.75in,top=0.6in,right=0.75in,bottom=0.6in]{geometry} % Document margins

\name{Leonardo V. Castorina} % Your name
\begin{document}
{\centering
    \href{https://github.com/universvm}{\faGithub \ universvm} \ \ \   
\href{https://www.linkedin.com/in/leonardo-castorina/}{\faLinkedin \ in/leonardo-castorina} \ \ \  
{\faPhone \ +44 7766 06 2571} \ \ \   
 \href{mailto:leonardo.castorina98@gmail.com}{\faEnvelope \ leonardo.castorina98@gmail.com}
\par}

%----------------------------------------------------------------------------------------
%	WORK EXPERIENCE SECTION
%----------------------------------------------------------------------------------------

\begin{rSection}{Experience}

\begin{rSubsection}{Osmitau Technologies}{Mar 2019 - Present}{– CEO \& Co-Founder}{Edinburgh (Scotland)}
\item Developing a hybrid drowsiness detection tool for the automotive industry.
\item Co-authored a white paper on automotive safety \fontsize{8pt}{12pt}\selectfont\href{https://osmitau.com/solutions.html}{ \faGlobe \ Read Here}
\end{rSubsection}

%------------------------------------------------

\begin{rSubsection}{IBM}{Jun 2019 - Sep 2019}{– Extreme Blue Technical Intern}{Hursley (UK)}
\item Built a search-planning interface for the search of vulnerable missing people with the UK police. 
\item Used LSTM, Convolutional Neural Networks (CNN), OpenStreetMap and statistics to derive an adaptive location radius which saved $\sim$1.5 hours per search.
\end{rSubsection}

%------------------------------------------------


\begin{rSubsection}{P\&G}{Jul 2018 - Sep 2018}{– R\&D Intern (Personal Healthcare)}{Egham (UK)}
\item Analysed the literature to verify the health benefits of a range of products. 
\item Planned a market and consumer research to identify the best strategy to enter the market. 
\item Designed a prototype app using Sketch to test customer behaviour.


\end{rSubsection}

%------------------------------------------------


\begin{rSubsection}{Swiss Institute of Bioinformatics}{May 2017 - Jul 2017}{– Data Visualization Intern}{Geneva (Switzerland)}
\item Redesigned ancestor graph for protein entries in \href{https://www.nextprot.org}{NeXtProt}. \fontsize{8pt}{12pt}\selectfont \href{https://github.com/universvm/nextprot-ancestor-graph}{ \faGithub \ View on Github}
\normalsize \item Developed Polymer.js components for
\href{https://glyconnect.expasy.org}{GlyConnect} to represent biological data.
\fontsize{8pt}{12pt}\selectfont
\href{https://github.com/universvm/glyco-balls}{ \faGithub \ View on Github}
\end{rSubsection}

%------------------------------------------------


\end{rSection}


\begin{rSection}{Projects}

\begin{rSubsection}{Dissertation: Explainable Deep Learning for {\it de novo} Protein Design (1st Class)}{}{}{}
\item Used 3D CNN to predict the identity of the residues required to obtain a protein shape. 
\item Explored explainability using activation maps and visualized accuracy and precision with Pymol.  
\end{rSubsection}


%------------------------------------------------

\begin{rSubsection}{PredictTheBurgh (Hackathon)}{\fontsize{8pt}{12pt}\selectfont \href
{https://github.com/universvm/PredictTheBurgh}{ \faGithub \ View on Github}}{}{}
\item Applied LSTM, GRU and CNN to predict stock prices and exposed flaws in published papers.
\end{rSubsection}

%------------------------------------------------

\begin{rSubsection}{Class Representative}{}{}{}
\item Visualized textual and numerical feedback using Seaborn and Pandas to suggest areas of improvement for the course organizers.
\end{rSubsection}

%------------------------------------------------
\end{rSection}

%----------------------------------------------------------------------------------------
%	EDUCATION SECTION
%----------------------------------------------------------------------------------------

\begin{rSection}{Education}

{\bf The University of Edinburgh} \hfill {\em Sep 2020 - Present} \\ 
MScR. Biomedical Artificial Intelligence \\

{\bf The University of Edinburgh} \hfill {\em Sep 2016 - May 2020} \\ 
BSc. (Hons) Biochemistry \\

\end{rSection}
%----------------------------------------------------------------------------------------
%	TECHNICAL STRENGTHS SECTION
%----------------------------------------------------------------------------------------

\begin{rSection}{Skills}

\begin{tabular}{ @{} >{\bfseries}l @{\hspace{3ex}} l }
Programming & Python, JavaScript, CSS \& HTML \\
Tech & TensorFlow, Keras, Sci-Kit Learn, NumPy, Pandas, Seaborn, Jupyter, Docker, Git \\
Software & Adobe Photoshop, Adobe Illustrator, Sketch
\end{tabular}

\end{rSection}

%----------------------------------------------------------------------------------------
%	Languages SECTION
%----------------------------------------------------------------------------------------

\begin{rSection}{Languages}

English (Fluent), Italian (Fluent), French (DELFB2) 
\end{rSection}

\fontsize{8pt}{12pt}\selectfont
\begin{flushright}
Last updated: \today \ \  \href{https://github.com/universvm/cv}{Latest: https://github.com/universvm/cv}
\end{flushright}

\end{document}
